% ---------------------------------------------------------------------------
%  PART II -- Objective question bank
% ---------------------------------------------------------------------------
\chapter{Objective Question Bank with Answers}
\chaptermark{Objective Question Bank}

Every multiple-choice question, fill-in-the-blank and one-mark short answer
recorded for HS248AT --- \textbf{76 items} in three sections. The
multiple-choice section is arranged \emph{by paper}, because the objective
part of a paper is written against the clock and it helps to know what one
paper's worth of them looks like.

\section{A. Multiple Choice Questions}

\subsection*{A.1\enspace From the CIE quiz papers (with official answers)}
\addcontentsline{toc}{section}{A.1 From the CIE quiz papers}

\begin{objtable}
\qn{1} & What is the fundamental principle of harmony in nature and existence?
  \optsinline{Competition}{Coexistence}{Dominance}{Isolation} &
  \textbf{(b) Coexistence} \\ \hline
\qn{2} & What is the role of human beings in maintaining harmony in nature?
  \optsinline{Exploiting natural resources}{Preserving and nurturing nature}%
  {Dominating other species}{Isolating from the environment} &
  \textbf{(b) Preserving and nurturing nature} \\ \hline
\qn{3} & Which of the following best describes the Physical Order?
  \optsinline{Includes only human-made objects}{Comprises natural elements
  like soil, water, air and minerals}{Refers exclusively to living beings}%
  {Focuses on mental and spiritual aspects} &
  \textbf{(b) Comprises natural elements like soil, water, air and minerals}
  \\ \hline
\qn{4} & What distinguishes the Human Order from the other three orders?
  \optsinline{Physical strength}{Intelligence and moral values}{Ability to
  photosynthesise}{Instinctual behaviour only} &
  \textbf{(b) Intelligence and moral values} \\ \hline
\qn{5} & In the context of coexistence, what is the meaning of `prosperity'?
  \optsinline{Material wealth accumulation}{Holistic well-being and
  happiness}{Physical strength and power}{Technological superiority} &
  \textbf{(b) Holistic well-being and happiness} \\ \hline
\qn{6} & Comprehensive Human Goals include all the following except:
  \optsinline{Resolution}{Prosperity}{Fearlessness (trust)}{Isolation} &
  \textbf{(d) Isolation} \\ \hline
\qn{7} & Gratitude is considered a \blank\ value in relationships.
  \optsinline{Cultural}{Financial}{Universal}{Personal} &
  \textbf{(c) Universal} \\ \hline
\qn{8} & The word ``society'' is primarily used in the context of human
  \blank\ relationship.
  \optsinline{Human}{Nature}{Both}{None} &
  \textbf{(a) Human} \\ \hline
\qn{9} & Prosperity deals with:
  \optsinline{Right understanding in the self}{Fulfilment in relationship}%
  {Ensuring more than required physical facility}{None} &
  \textbf{(c) Ensuring more than required physical facility} \\ \hline
\qn{10} & \blank\ is the feeling of responsibility towards the body of my
  relative.
  \optsinline{Care}{Guidance}{Respect}{Affection} &
  \textbf{(a) Care} (mamata) \\ \hline
\qn{11} & \blank\ means harmony within myself.
  \optsinline{Pleasure}{Happiness}{Excitement}{Fearlessness} &
  \textbf{(b) Happiness} \\ \hline
\qn{12} & The fourth order of nature is \blank.
  \optsinline{Bio order}{Animal order}{Human order}{Material order} &
  \textbf{(c) Human order} \\ \hline
\qn{13} & The continuity of a plant species is maintained in nature by the
  \blank\ method.
  \optsinline{Constitution conformance}{Seed conformance}{Right value \,/\,
  sanskaar}{Breed conformance} &
  \textbf{(b) Seed conformance} \\ \hline
\qn{14} & There is \blank\ among the four orders of nature.
  \optsinline{Inherent struggle}{No dependency}{Mutual fulfilment}%
  {Opposition} &
  \textbf{(c) Mutual fulfilment} \\ \hline
\qn{15} & Our natural acceptance is to live in \blank\ with nature.
  \optsinline{Harmony}{Struggle for survival}{Opposition}{None of the above} &
  \textbf{(a) Harmony} \\ \hline
\qn{16} & Production and work for physical facilities leads to \blank\ in
  family and \blank\ with nature.
  \optsinline{Prosperity, existence}{Happiness, existence}{Happiness,
  co-existence}{Prosperity, co-existence} &
  \textbf{(d) Prosperity, co-existence} \\ \hline
\end{objtable}

\subsection*{A.2\enspace CIE-II --- 7 May 2025}
\addcontentsline{toc}{section}{A.2 CIE-II, 7 May 2025}

\begin{objtable}
\qn{17} & Which of the following is considered the foundation of a strong
  human relationship?
  \optsinline{Intelligence}{Trust}{Wealth}{Appearance} &
  \textbf{(b) Trust} --- the \emph{adhara mulya}, foundation value
  \seeq{31}\seeq{32} \\ \hline
\qn{18} & Disrespect Arising out of
  \begin{optlist}
  \item Lack of communication
  \item Lack of knowledge
  \item Differentiation leading to Discrimination
  \item Lack of Wealth
  \end{optlist} &
  \textbf{(c) Differentiation leading to Discrimination} \seeq{36} \\ \hline
\qn{19} & Care is the feeling
  \begin{optlist}
  \item of responsibility and commitment for nurturing and protection of the
        Body of my relative.
  \item of responsibility and providing facility for nurturing and protection
        of the Body of my relative
  \item of responsibility of making the other feel comfortable with respect to
        Body of my relative
  \item of showing affection to the Body of my relative
  \end{optlist} &
  \textbf{(a)} --- the text defines care (\emph{mamata}) as the feeling to
  \emph{nurture and protect the body of our relative}, taken up as a
  responsibility \seeq{39} \\ \hline
\qn{20} & Fulfilment of relationship means
  \begin{optlist}
  \item Making effort for self development, i.e.\ development of one's own
        competence
  \item Making effort for mutual development, i.e.\ development of one's own
        competence and being of help to the other in developing their
        competence
  \item Making effort for mutual understanding, i.e.\ understanding of one's
        own competence and being of help to the other in understanding their
        competence
  \item Making effort for mutual in acquiring prosperity
  \end{optlist} &
  \textbf{(b) Mutual development} --- the response mode is defined as
  ``a sense of responsibility to improve our own competence and the other's
  competence; we work for mutual fulfilment'' \seeq{48} \\ \hline
\qn{21} & With the complete understanding of respect
  \begin{optlist}
  \item \ldots we all are the same in terms of intelligent, prosperity and
        status.
  \item \ldots we all are the same in terms of Knowledge, family, society.
  \item \ldots we all are the same in terms of competence, prosperity and
        status.
  \item \ldots we all are the same in terms of intention, program and
        potential.
  \end{optlist} &
  \textbf{(d) intention, program and potential} \seeq{35} \\ \hline
\end{objtable}

\begin{note}
\textbf{\sffamily A note on item 21.} The text's own triad is \emph{basic
aspiration} (we both want continuous happiness and prosperity),
\emph{programme of action} (to understand and live in harmony at all four
levels) and \emph{potential} (the activities and powers of the Self are the
same in both). The paper writes ``intention'' where the text writes
``aspiration\,/\,desire'', but option (d) is the only one carrying the triad
at all --- and note that (c) is a deliberate trap: \textbf{competence} is
precisely the one thing in which we are \emph{not} the same.
\end{note}

\subsection*{A.3\enspace Improvement CIE --- June 2025}
\addcontentsline{toc}{section}{A.3 Improvement CIE, June 2025}

\begin{objtable}
\qn{22} & The participation of the human being in ensuring the role of
  physical facility in nurture, protection and providing means for the body is
  called its \blank
  \optsinline{Utility value}{Artistic value}{Activity}{Energy} &
  \textbf{(a) Utility value} \seeq{4} \\ \hline
\qn{23} & When something is active or has activity, it is called as \blank
  \optsinline{space}{unit}{energy}{value} &
  \textbf{(b) unit} --- ``only units are active; when something is active or
  has activity, we call it a unit'' \seeq{65} \\ \hline
\qn{24} & The following are the characteristics of space except \blank
  \optsinline{Unlimited}{No activity}{Energy in equilibrium}{Self-organized} &
  \textbf{(d) Self-organized} --- space is \emph{not} self-organised; it is
  that in which \emph{self-organisation is available} \seeq{61}\seeq{65}
  \\ \hline
\qn{25} & The natural characteristic of material order is \blank
  \optsinline{Composition}{Decomposition}{Composition\slash Decomposition}%
  {Nurture\slash worsen} &
  \textbf{(c) Composition\slash Decomposition} \seeq{52}\seeq{54} \\ \hline
\qn{26} & Conformance of material order is named as \blank
  \optsinline{constitution conformance}{Seed conformance}{Breed conformance}%
  {sanskaar conformance} &
  \textbf{(a) constitution conformance} \seeq{55} \\ \hline
\end{objtable}

\subsection*{A.4\enspace Semester End Examination --- June/July 2025}
\addcontentsline{toc}{section}{A.4 SEE, June/July 2025}

\begin{objtable}
\qn{27} & How does ensuring justice in relationships contribute to society?
  \begin{optlist}
  \item It leads to fearlessness and co-existence
  \item It leads to prosperity and right understanding
  \item It promotes self-regulation and health
  \item It ensures better exchange and storage
  \end{optlist} &
  \textbf{(a) fearlessness and co-existence} --- \emph{Nyaya-Suraksha} is the
  dimension that leads to \emph{abhaya} and \emph{saha-astitva}
  \seeq{44}\seeq{69} \\ \hline
\qn{28} & Imagination is continuous with
  \optsinline{Time}{Soul}{Body}{None} &
  \textbf{(a) Time} --- ``imaging continues with time; it is taking place all
  the time'' \seeq{24} \\ \hline
\qn{29} & Identify the solution which helps human being to transform from
  animal consciousness to human consciousness.
  \optsinline{Right understanding}{Realization}{Value education}{Physical
  facilities} &
  \textbf{(a) Right understanding} --- the transformation is accomplished
  ``only by working for right understanding as the first priority''
  \seeq{22} \\ \hline
\qn{30} & Which of the following is a fundamental human aspiration related to
  relationships?
  \begin{optlist}
  \item Physical comfort.
  \item Financial security.
  \item Mutual understanding and fulfillment.
  \item Accumulating wealth.
  \end{optlist} &
  \textbf{(c) Mutual understanding and fulfillment} \seeq{19} \\ \hline
\qn{31} & Which of the following capacity leads to desires
  \optsinline{Power}{Expectation}{Realization}{Thoughts} &
  \textbf{(c) Realization} --- see the note below \seeq{24} \\ \hline
\qn{32} & Which of the following is considered the foundation of a strong
  human relationship?
  \optsinline{Intelligence}{Trust}{Wealth}{Appearance} &
  \textbf{(b) Trust} \seeq{31}\seeq{32} \\ \hline
\qn{33} & Respect in a relationship means:
  \begin{optlist}
  \item Agreeing to everything
  \item Controlling the other person's actions
  \item Valuing the other person's opinions and boundaries
  \item Always trying to win argument.
  \end{optlist} &
  \textbf{(c) Valuing the other person's opinions and boundaries} --- the
  option nearest to \emph{right evaluation} \seeq{35} \\ \hline
\qn{34} & The participation of the human being in ensuring the role of
  physical facility to help and preserve its utility is called its \blank
  \optsinline{Utility value}{Artistic value}{Activity}{Energy} &
  \textbf{(a) Utility value} \seeq{4} \\ \hline
\qn{35} & A harmonious world is created by values at 4 levels. These are:
  \begin{optlist}
  \item Home, family, society, country
  \item Individual, family, society, universe
  \item School, home, office, temple
  \item None of these
  \end{optlist} &
  \textbf{(b) Individual, family, society, universe} \seeq{5} \\ \hline
\qn{36} & When nature is submerged in space it is known as
  \optsinline{Conformance}{Acceptance}{Mixing}{Co-existence} &
  \textbf{(d) Co-existence} --- see the note below \seeq{62} \\ \hline
\end{objtable}

\begin{note}
\textbf{\sffamily Three items on this paper need care.}

\textbf{Item 31 --- ``which capacity leads to desires''.} The text gives the
self-organised sequence as \textbf{Realization \& Understanding (1, 2)
$\rightarrow$ Desire (3) $\rightarrow$ Thought (4) $\rightarrow$ Expectation
(5)}. On that sequence the capacity that leads to desire is
\textbf{Realization}, and that is the answer to give. Be aware of the trap:
the text \emph{also} describes today's enslaved (\emph{partantra}) flow, in
which the direction is reversed --- sensation sets thoughts, and thoughts set
desires. `Thoughts' is therefore true of how we live now, but `Realization' is
the capacity the text names.

\textbf{Item 35 --- the four levels.} The text's four levels of living are
\emph{individual, family, society, nature\,/\,existence}. The paper's option
(b) writes ``universe'' for nature\,/\,existence; it is nonetheless the only
option carrying the right sequence.

\textbf{Item 36 --- nature submerged in space.} The text's own equation is
\textbf{Existence $=$ Nature submerged in Space}, so the strictly correct word
is \emph{existence}, which is not on offer. Since the text immediately adds
that \textbf{existence is co-existence}, option (d) is the intended answer.
\end{note}

\section{B. Fill in the Blanks and One-Liners}

\begin{objtable}
\qn{1} & \blank\ and Respect are considered the foundational values of a
  relationship. & \textbf{Trust} (vishvas) \\ \hline
\qn{2} & Visualising a universal harmonious order in society includes the idea
  of an \blank\ Society and Universal Order, from family to world family. &
  \textbf{Undivided} (Akhanda Samaja) \\ \hline
\qn{3} & Prosperity and fearlessness (trust) are part of the comprehensive
  human goals in \blank\ in society. &
  \textbf{Right understanding} (samadhana) \\ \hline
\qn{4} & There are \blank\ comprehensive human goals. &
  \textbf{Four} --- right understanding, prosperity, fearlessness (trust),
  co-existence \\ \hline
\qn{5} & The human goal at the level of the individual is \blank. &
  \textbf{Right understanding} (samadhana) \\ \hline
\qn{6} & The human goal at the level of the family is \blank. &
  \textbf{Prosperity} (samriddhi) \\ \hline
\qn{7} & The human goal at the level of society is \blank. &
  \textbf{Fearlessness \,/\, trust} (abhaya) \\ \hline
\qn{8} & The human goal at the level of nature is \blank. &
  \textbf{Co-existence} (saha-astitva) \\ \hline
\qn{9} & Human--human interaction (its recognition, fulfilment and evaluation
  leading to mutual happiness) is called \blank. &
  \textbf{Justice} (nyaya) \\ \hline
\qn{10} & Human--rest-of-nature relation (recognition, fulfilment, evaluation
  leading to mutual prosperity) is called \blank. &
  \textbf{Preservation} (suraksha) --- enrichment, protection, right
  utilisation \\ \hline
\qn{11} & Human values are essential for \blank. &
  \textbf{Harmony at all four levels of living} --- and hence for mutual
  fulfilment and mutual happiness in relationship \\ \hline
\qn{12} & Mutual happiness in human--human interaction leads to \blank. &
  \textbf{Justice} (nyaya) --- mutual happiness ensured is the fourth element
  of justice \\ \hline
\qn{13} & \blank\ is the foundation value (adhara mulya) in relationship. &
  \textbf{Trust} (vishvas) \\ \hline
\qn{14} & What do you understand by complete value? &
  \textbf{Love} (prema) --- the feeling of being related to all; hence called
  \emph{purna mulya} \\ \hline
\qn{15} & The value or participation of different orders in existence is also
  referred to as \blank. &
  \textbf{Natural characteristic} (svabhava) \\ \hline
\qn{16} & \emph{Sah-astitva} means \blank. & \textbf{Co-existence} \\ \hline
\qn{17} & Which activity of the self indicates that ``there is
  complementarity in nature and no opposition''? &
  \textbf{Realization} (and understanding) --- activities 1 and 2 of the Self
  \\ \hline
\qn{18} & Respect means \blank. &
  \textbf{Right evaluation} --- to be evaluated as I am \\ \hline
\qn{19} & The three mistakes in evaluation are \blank, \blank\ and \blank. &
  \textbf{Over evaluation} (adhi-mulyana), \textbf{under evaluation}
  (ava-mulyana), \textbf{otherwise evaluation} (a-mulyana) \\ \hline
\qn{20} & Ensuring justice in relationship, on the basis of values, leads to
  \blank\ in society. &
  \textbf{Fearlessness \,/\, trust} (abhaya) \\ \hline
\qn{21} & The two mechanisms of self-exploration are \blank\ and \blank. &
  \textbf{Natural acceptance} and \textbf{experiential validation} \\ \hline
\qn{22} & Realization and understanding are tested for \blank, \blank\ and
  \blank. &
  \textbf{Assurance, satisfaction and universality} (same for all time, space
  and individual) \\ \hline
\qn{23} & The needs of the Self are \blank\ and \blank; the needs of the Body
  are \blank\ and \blank. &
  \emph{Self:} \textbf{continuous and qualitative}. \emph{Body:}
  \textbf{temporary and quantitative}. \\ \hline
\qn{24} & Sanyama is the basis of \blank. & \textbf{Svasthya} (health)
  \\ \hline
\qn{25} & The five dimensions of human endeavour are \blank. &
  \textbf{Siksha-Sanskara, Svasthya-Sanyama, Nyaya-Suraksha, Utpadana-Karya,
  Vinimaya-Kosha} \\ \hline
\qn{26} & The three aspects of suraksha are \blank. &
  \textbf{Enrichment} (samvardhana), \textbf{Protection} (sanrakshana),
  \textbf{Right utilisation} (sadupayoga) \\ \hline
\qn{27} & Production must be through a \blank\ process, in harmony with
  nature. &
  \textbf{Cyclical} (avartansheel) process --- cyclic, and enriching every
  unit \\ \hline
\qn{28} & Existence $=$ \blank\ $+$ \blank. &
  \textbf{Space $+$ Units} (in space); i.e.\ \textbf{nature submerged in
  space} \\ \hline
\qn{29} & Space is \blank\ energy; units are \blank. &
  Space is \textbf{energy in equilibrium} (constant energy); units are
  \textbf{energised and active} \\ \hline
\qn{30} & Conformance in the four orders is \blank, \blank, \blank\ and
  \blank. &
  \textbf{Constitution, Seed, Breed, Sanskar} conformance \\ \hline
\qn{31} & The innateness of the human order at the level of `I' is \blank. &
  \textbf{The will to live with happiness} \\ \hline
\qn{32} & The svabhava of the self (`I') in human beings is \blank. &
  \textbf{Perseverance} (dhirata), \textbf{bravery} (virata) and
  \textbf{generosity} (udarata) \\ \hline
\qn{33} & SVDD, SSDD and SSSS stand for \blank. &
  \textbf{Sadhan Viheen Dukhi Daridra; Sadhan Sampann Dukhi Daridra; Sadhan
  Sampann Sukhi Samriddha} \\ \hline
\qn{34} & Svatva, Swatantrata and Swarajya mean \blank. &
  \textbf{Innateness; self-organisation} (harmony in oneself);
  \textbf{self-expression} (harmony with others) \\ \hline
\qn{35} & The four elements of justice are \blank. &
  \textbf{Recognition of values, fulfilment, evaluation, mutual happiness
  ensured} \\ \hline
\end{objtable}

\section{C. One-Mark Short Answers}

The Improvement CIE of \textbf{18 June 2026} broke with the earlier pattern:
its Part~A was not multiple choice but five one-mark short-answer questions,
to be written in ten minutes. There is nothing to guess at, so these are worth
rehearsing as written answers.

\begin{objtable}
\qn{1} & Give any two definitions for the term co-existence. \mmk{1} &
  Any two of: \textbf{(i)} to exist together and in mutual tolerance;
  \textbf{(ii)} to learn to recognise and live with difference;
  \textbf{(iii)} a relationship in which none of the parties is trying to
  destroy the other \seeq{63} \\ \hline
\qn{2} & State the four orders of nature. \mmk{1} &
  \textbf{Material} (padartha avastha), \textbf{Pranic\,/\,Bio} (prana
  avastha), \textbf{Animal} (jiva avastha), \textbf{Human\,/\,Knowledge}
  (gyana avastha) \seeq{51} \\ \hline
\qn{3} & Identify the svabhava\,/\,value of the self (`I') in human beings.
  \mmk{1} &
  \textbf{Perseverance} (dhirata), \textbf{bravery} (virata) and
  \textbf{generosity} (udarata) \seeq{54} \\ \hline
\qn{4} & Summarize 2 things that human beings have to do to realize the mutual
  fulfilment in nature. \mmk{1} &
  \textbf{(i) Understand} the harmony and interconnectedness among the four
  orders --- and our own needs --- correctly; \textbf{(ii) live and produce
  accordingly}, through a cyclic (avartansheel) process that enriches every
  unit \seeq{59}\seeq{60} \\ \hline
\qn{5} & Differentiate between Units and Space. \mmk{1} &
  \textbf{Units} are limited, active, energised, recognise and fulfil their
  relation with other units, and are self-organised. \textbf{Space} is
  unlimited and all-pervading, no-activity, energy in equilibrium, reflecting
  and transparent, and self-organisation is available in it \seeq{65}
  \\ \hline
\end{objtable}

\begin{note}
\textbf{\sffamily On item 4.} The text does not give a numbered list of ``two
things'' under that heading. What it does say, in diagnosing why the human
order alone fails to be mutually fulfilling, is that \emph{``we have not
understood the harmony that exists between these orders. We have not even
understood our own needs properly, nor have we understood harmonious ways to
fulfil our needs.''} The two things to do are therefore the two failures put
right: \textbf{understand}, and then \textbf{live and produce accordingly}.
That is the answer given above; it is derived from the text rather than quoted
from a list in it.
\end{note}
